\begin{titlepage}
\color{navy}\sffamily
\noindent\large The Energy Balance Project
\vspace{38pt}
\par\noindent\fontsize{28}{34}\selectfont\bfseries What an Economy\par Must Keep Alive
\vspace{24pt}
\par\noindent\fontsize{17}{23}\selectfont The physical theory and design of an Energy Balance Unit
\vspace{27pt}
\par\noindent\normalfont A continuous explanatory edition\par\small Expanded theoretical integration
\vfill
\noindent\sffamily\Large Konrad Grzyb
\par\normalsize Author and project originator
\vspace{20pt}
\par\small AUTHOR-REVIEW MANUSCRIPT\par October 2026
\end{titlepage}

\chapter*{Preface}
Every generation inherits a world that makes its activity possible. Money helps
people exchange, coordinate and plan within that world. Yet a completed sale
cannot, by itself, tell us whether the water line will still work, whether a
field will still produce food or whether a community can absorb its next
disturbance. Those are physical questions. This book asks what an economy would
look like if its account of action had to answer to them as well.

We begin with the ordinary reasons money arose, and with the incentives it now
creates. From there the argument moves through light, energy, thermodynamics
and the living body's continuous work of homeostasis. The point of these
comparisons is not to pretend that an economy is a photon or an organism. It
is to look for a disciplined way to connect local actions to the conditions
that keep a larger system functioning.

The proposed Energy Balance Unit enters only after that problem is clear.
We declare what physical state is represented, how an action changes it, and
what potential expresses a chosen functional comparison. The exact finite
EBU is the decrease of that potential across the complete accepted change.
An integral along the change shows why an initial slope is not enough and why
the finite result can be calculated without asking one actor to know the
state of the whole planet.

The later chapters ask what happens when actions share sources, when physical
effects cross several steps, and when the field itself changes. They also ask
how a future EBU system could keep its physical permissions, measurements,
attributions and actor decisions distinct. An equation may be exact while the
world chosen for it remains a hypothesis. We therefore explain the mathematics
fully, then show what an implementation and an independent test would have to
establish before anyone could claim more.

The physical consequence, its mathematical value, a personal attribution and
an institution's rule for using that attribution will remain distinguishable
throughout. Keeping them connected without making them identical is the
discipline on which EBU's promise depends. We can begin, then, where people
already make their choices: with money, the work it does for us, and the
physical conditions it cannot describe by itself.

\section*{How to read the extended argument}
The opening chapters build the physical vocabulary without assuming advanced mathematics. The central sequence then moves from paths and finite changes to coalition tables, recursive interaction and exact discrete Taylor structure. It asks how a nonlinear response or a generator supplies the endpoints before adding the narrower probability and entropy interpretations. Each major result is followed by what it could make possible for an EBU system and what remains a matter of evidence or institutional choice.

The final chapters return to routes, attribution, measurement, implementation and access. The appendices contain the longer calculus and generator proofs, statistical counterexamples and a worked claim review. They can be read after the main argument without breaking its continuity.

A definition declares an object. A theorem or corollary derives a consequence under stated assumptions. Constructed examples teach the mathematics. Physical interpretations require their named models, while hypotheses and future studies remain open. This edition integrates independently cleared source mathematics; it remains an author-review book requiring a separate independent book audit.
